% Method/settings fragment for the team's CVPR report; not a complete paper.
% This fragment describes the delivered implementation and contains no performance claims.
\section{Method}
\label{sec:method}

We study binary detection of AI-generated images under a strict prohibition on
pretrained weights. The target is $y=1$ for generated images and $y=0$ for real
images. All networks start from random initialization. We investigate local RGB
appearance, a fixed residual input, and prediction consistency under JPEG
re-encoding through controlled ablations.

\paragraph{Aligned image pairs.}
We correct EXIF orientation and convert each image to RGB. For local inputs, we
preserve the original resolution and apply symmetric reflection padding when a
side is shorter than 224 pixels. During training, an image is horizontally
flipped with probability $0.5$ and re-encoded as JPEG with integer quality sampled
uniformly from $[60,100]$, with chroma subsampling disabled. Re-encoding is applied
to the entire transformed image before cropping. The normal and JPEG views use
the same random $224\times224$ crop. Development inference uses a deterministic
center crop of the normal view. The global baseline instead resizes the whole
image to $224\times224$ using bilinear interpolation.

\paragraph{RGB and auxiliary encoders.}
The RGB encoder uses a ResNet-18-style structure with a $3\times3$, stride-one
stem, no initial max pooling, and two residual blocks in each of four stages.
The stage widths are $32,64,128,256$. Global average pooling yields a
256-dimensional feature. RGB values in $[0,1]$ are normalized using channel
statistics estimated exclusively from the 5,600 training images.

For the residual variant, we compute a fixed, channel-wise high-pass input
\begin{equation}
 r(x)=x-\operatorname{AvgPool}_{3\times3}(x),
\end{equation}
using reflection padding and stride one. This operation is applied to raw RGB,
and its output is divided by the same training-channel standard deviations;
we do not subtract the RGB means from the residual. The auxiliary encoder has
three $3\times3$, stride-two convolution--BatchNorm--ReLU stages with widths
$32,64,128$, followed by global average pooling. We concatenate its 128-dimensional
output with the RGB representation and predict one logit using a linear head.
The single RGB model has 2,795,297 trainable parameters; each fused model has
2,888,897. The capacity control uses the identical auxiliary architecture and
head but supplies normalized RGB instead of the fixed residual. It therefore
matches the residual variant's trainable parameter count exactly.

\paragraph{Supervision and consistency.}
Let $z=f_\theta(x)$ and $\tilde z=f_\theta(\tilde x)$ denote the logits of an
aligned normal/JPEG pair. Both views are concatenated into a single forward
pass, so BatchNorm observes the same mixed batch. All experimental groups use
the same paired supervision,
\begin{equation}
 \mathcal L_{\mathrm{sup}}=\tfrac12\big[
 \operatorname{BCE}(z,y)+\operatorname{BCE}(\tilde z,y)\big].
\end{equation}
For $P=\operatorname{Bernoulli}(\sigma(z))$ and
$Q=\operatorname{Bernoulli}(\sigma(\tilde z))$, we use
\begin{align}
 M&=\tfrac12(P+Q),\\
 \operatorname{JS}(P,Q)&=\tfrac12\operatorname{KL}(P\Vert M)
                       +\tfrac12\operatorname{KL}(Q\Vert M),\\
 \mathcal L&=\mathcal L_{\mathrm{sup}}+\lambda_t\operatorname{JS}(P,Q).
\end{align}
The implementation computes Bernoulli log probabilities with log-sigmoid and
the mixture with log-add-exp in float32. For a total of $T$ optimization steps,
the enabled coefficient is
$\lambda_t=0.1\min\{1,t/\lceil0.1T\rceil\}$ for one-indexed step $t$.
The disabled setting uses $\lambda_t=0$ while retaining both views and their BCE
terms. At the frozen batch size there are 350 steps per epoch and 7,000 in total;
the ramp reaches $0.1$ at step 700.

\section{Experimental Settings}
\label{sec:settings}

The fixed labeled-data split contains 5,600 training, 1,200 development, and
1,200 held-out images, with equal class counts in each split. The official
unlabeled test set is excluded from training, normalization, and model
selection. The training-channel statistics are estimated after resizing each
training image to $128\times128$ with equal image weighting. The means are
$(0.483540,0.458205,0.404792)$ and the standard deviations are
$(0.270488,0.262578,0.276671)$.

Every research group uses 20 epochs, batch size 16 image pairs, AdamW with
initial learning rate $10^{-3}$ and weight decay $10^{-4}$, cosine annealing,
automatic mixed precision, and seeds 17, 29, and 43. A batch of 16 pairs is
forwarded as 32 view tensors; it is not a batch of 16 total views. We select
checkpoints using development accuracy at threshold $0.5$, breaking ties by
lower BCE and then earlier epoch. Reported development results are therefore
selection-set measurements, not independent generalization estimates.

The local factorial design compares A (RGB), B (RGB plus residual), C (RGB plus
consistency), and D (both). E and F replace the residual branch with a second
RGB branch, without and with consistency, respectively; these capacity controls
must be available before attributing an improvement specifically to residual
inputs. G measures the global-resize baseline. We report seed means and sample
standard deviations, as well as paired differences for matching seeds.

Real images in this dataset are all JPEG and generated images are all PNG;
813 labeled images have a short side below 224 pixels. Consequently, format,
source, resolution, and padding artifacts are potential shortcuts. Random
splits and paired JPEG augmentation alone do not establish cross-source
generalization or robustness. The independent evaluation, capacity controls,
and frozen held-out confirmation are reported separately by the evaluation
owner. Full-data refitting occurs only after the final method and recipe are
frozen and does not use development or held-out results to select an epoch.
